\documentclass[11pt]{book}

\usepackage{graphicx}
\usepackage{eso-pic}
\usepackage{tcolorbox}
\usepackage{polyglossia}
\setdefaultlanguage[indentfirst=false,forceheadingpunctuation=false]{russian}
\setotherlanguages{english}

\RequirePackage[a4paper, headsep = 0.5 \headsep, left=2.5cm, right=2.1cm, top=2cm, bottom=2.1cm]{geometry}

\setmainfont{Times New Roman}
\newfontfamily\cyrillicfont{Times New Roman}[Script=Cyrillic, Ligatures=TeX]

\usepackage{enumitem}
\setlist[itemize, 1]{label =\raisebox{-0.3\height}{\scalebox{1.8}{\textbullet}}}

\title{Феноменология: Основы}
\author{Дэн Захави}
\date{2019}
\begin{document}

\AddToShipoutPictureBG*{\includegraphics[width=\paperwidth]{title.png}}

\null
    \thispagestyle{empty}%
    \addtocounter{page}{-1}%
    \newpage

\tableofcontents

\chapter{Введение}

Феноменология считается одной из доминирующих традиций в философии XX века. Ее основал Эдмунд Гуссерль (1859-1938)\footnote{Эдмунд Гуссерль пришел в философию в относительно позднем возрасте. Он родился в Проснице, Моравия (тогда входившая в состав Австрийской Империи), изучал физику, математику, астрономию и философию в Лейциге, Берлине и Вене. В 1883г он защитил докторскую диссертацию по математике, и в последующие годы серьезно заинтересовался философией , отчасти благорадя посещению лекций психолога и философа Франца Брентано. Первая крупная работа Гуссерля --- ''Логические исследования'' --- была опубликована в 1900-1901гг, и именно на ее основе он получил приглашение в Геттинген, где преподавал с 1901 по 1916 год. Следующая крупная работа --- ''Идеи к чистой феноменологии и феноменологической философии I'' --- была опубликована в 1913г. В 1916г Гуссерль переезжает во Фрайбург, где получает кафедру философии, прежде занимаемую Генрихом Риккертом. В последующие годы его ассистентами работают Эдит Штайн и Мартин Хайдеггер, который, когда Гуссерль выйдет на пенсию в 1928г, займет его место. Вскоре после выхода на пенсию Гуссерль публекует ''Формальную и трансцендентальную логику'' (1929) и ''Картезианские размышления'' (1931). В тридцатые годы Гуссерль страдал от условий, навязанных немецким национал-социалистическим режимом: лишенный какой-либо официальной академической деятельности из-за своего еврейского происхождения, он теряет право преподавать, публиковаться, и, в конечном итоге, свое немецкое гражданство. Он, однако, продолжает свою работу, в такие темные времена еще более настаивая на актуальности философии. В 1935г Гуссерля приглашают читать лекции в Вене и Праге, эти лекции в последствии лягут в основу его последней работы ''Кризис европейских наук и трансцендентальная феноменология'' (1936). Вскоре после его смерти 27 апреля 1938г молодому монаху-францисканцу Герману Лео ван Бреде удается котрабандой вывезти многочисленные исследовательские рукописи Гуссерля из Германии в монастырь в Бельгии. Перед началом Второй Мировой Войны в Левене был создан Архив Гуссерля, где до сих пор можно найти оригинальные рукописи.}, другие влиятельные сторонники --- Мартин Хайдеггер (1889-1976)\footnote{Хайдеггер начал изучать католическую теологию и средневековую философию во Фрайбурге, но в 1911г решил сосредоточиться на философии. В 1913г он защищает диссертацию, а два года спустя --- хабилитацию ''Учение Дунса Скотта о категориях и значении'', которая затем будет передана Риккерту --- философу, преемником которого был Гуссерль. Хайдеггер работает ассистентом Гуссерля в 1918-1923гг, и затем становится профессором университета в Марбурге. В 1927 году он публикует главную свою работу --- ''Бытие и время'' --- а в 1928 занимает кафедру Гуссерля во Фрайбурге. В 1929г Хайдеггер прочел свою знаменитую вступительную лекцию ''Что такое метафизика?''. После прихода к власти Гитлера Хайдеггер был избран ректором Фрайбургского университета и стал членом нацистской партии. Однако менее чем через год он уходит с этого поста и постепенно отходит от университетской политики. Как показывает публикация в 2014г его так называемых ''Черных тетрадей'', флирт Хайдеггера с нацизмом не был ни кратковременным, ни поверхностным. До 1944г он регулярно читал лекции, но после окончания войны ему запретили преподавать из-за политических убеждений, а в 1946г он был лишен звания профессора. В 1949г его восстанавливают в звании почетного профессора, и с тех пор и до своей смерти о читает много лекций и публикует такие центральные работы как ''Язык'' (1950), ''Размышления о строительстве и обитании'' (1951) и ''Вопрос технологии'' (1953).}, Жан-Поль Сартр (1905-1980)\footnote{Сартр изучал философию в Высшей нормальной школе, где встретил целое поколение ведущих французских интеллектуалов, в том числе Симону де Бовуар, Раймона Арона, Мориса Мерло-Понти, Симону Вейль, Эммануэля Мунье, Жана Ипполита и Клода Леви-Стросса. Его отношения с коллегой-философом Бовуар станут легендарными. В период 1931-1945гг Сартр преподает в средних школах Гавра, Лана и Парижа. С философией Гуссерля и Хайдеггера его знакомят в начале 1930гг Арон и Левинас, а в 1933-1934гг он отправляется в Берлин изучать феноменологию. Во второй половине тридцатых годов Сартр публикует четыре книги о сознании: ''Трансцендентность эго'', в которой рассматривается структура самосознания (1936), ''Воображение'' (1936) и ''Воображаемое'' (1940), а также ''Очерк теории эмоций'' (1939). Все четыре эти работы свидетельствуют о его знакомстве с философией Гуссерля. С начачлом Второй Мировой Войны Сартр был призван в армию и в 1940г попал в плен к немцам. После освобождения в 1941г он становится активным участником сопротивления и в 1943г публикует ''Бытие и ничто'', ясно демонстрирующий влияние Хайдеггера. В 1945г Сартр основывает литературно-политическое обозрение Les Temps Moderne (''Новые времена''), которое редактирует в течение многих лет (периодически совместно с Мерло-Понти). После войны Сартр решает прекратить преподавание и посвящает себя литературной и редакционной работе. Сартр, таким образом, один из очень немногих влиятельных философов XX века, который не работал в университете. В послевоенный период политическая активность Сартра возрастает, возрастают и его симпатии к марксизму и восхищение СССР, хотя членом коммунистической партии он так и не станет. Поддержка им СССР остается неизменной до вторжения в Венгрию в 1956г. В 1960г Сартр публикует свое второе основное произведение ''Критика диалектического разума'', что свидетельствует о его политической и социальной активности. В 1964г Сартру присудят Нобелевскую премию по литературе, но он из принципа откажется от нее. Он останется политически активным до позднего возраста, выступая против войны Франции в Алжире, активно выступая против войны во Вьетнаме вместе с Бертраном Расселом, а также поддерживая протесты в Париже в мае 1968г. На его похороны после смерти в апреле 1980г приде около 50тыс. человек.}, Морис Мерло-Понти (1908-1961)\footnote{Мерло-Понти изучал философию в престижной Высшей нормальной школе в Париже и был первым иностранным посетителем недавно созданного Архива Гуссерля в Левене в апреле 1939г. Свою первую книгу ''Структура поведения'' он публикует в 1942г, а свою, пожалуй, главную работу --- ''Феноменология восприятия'' --- в 1945г. В 1949г он становится профессором детской психологии и педагогики в Сорбонне, а в 1952г будет избран на кафедру философии в Коллеж де Франс, став самым молодым из когда-либо назначенных на эту должность, и будет занимать ее до своей смерти в мае 1961г. После Второй Мировой Войны Мерло-Понти все активнее занимается политикой и публикует ряд книг с эссе, среди которых ''Гуманизм и террор'' (1947), ''Смысл и бессмыслица'' (1948) и ''Приключения диалектики'' (1955). Параллельно с этими политическими интересами он продолжает преподавать, и многие из его лекций в Сорбонне и Коллеж де Франс свидетельствуют об обширном интересе к другим дисциплинам, включая детскую психологию, структурную лингвистику, этнологию и психоанализ. В 1960г выходит еще один том ''Знаков'', состоящий из эссе, а в 1964г будет посмертно опубликован отрывок ''Видимое и невидимое'', которое многие считают вторым главным произведением Мерло-Понти.} и Эммануэль Левинас (1906-1995)\footnote{Эммануэль Левинас родился в Каунасе в Литве (тогда входившей в состав Российской Империи). После окончания средней школы он уезжает во Францию, где начинает изучать философию в Страсбурге. В 1928г он обучается у Гуссерля и Хайдеггера во Фрайбурге, и затем совместно с Г.Пайффером переведет ''Картезианские размышления'' Гуссерля на французский. С публикацией в 1931г его докторской диссертации ''Теория интуиции в феноменологии Гуссерля'' за ним устанавливается репутация как одного из ведущих специалистов по немецкой феноменологии во Франции. К моменту начала Второй Мировой Войны Левинас стал гражданином Франции и был призван в армию. Как и Сартр, он вскоре попал в плен к немцам, где провел остаток войны, благодаря чему избежал участи большей части своей еврейской семьи, убитой немцами. После войны Левинас становится директором École Normale Israélite Orientale, а затем получает должности в университетах Пуатье (1961), Нантера (1967), а с 1973г --- в Сорбонне в Париже. Ранние послевоенные работы Левинаса ''От существования к существующему'' (1947), ''Время и другой'' (1948) и ''Открытие существования с Гуссерлем и Хайдеггером'' (1949) все еще демонстрируют его долг как перед Гуссерлем, так и перед Хайдеггером, но уже рассматривают характерные для собственного мышления Левинаса темы, такие как роль другого и отношение между этикой и онтологией. Рассмотрение этих тем достигает своей кульминации в ''Тотальность и бесконечность'' (1961) и ''Иное, чем бытие или за пределами сущности'' (1974). Помимо множества философских книг, Левинас опубликует также серию талмудических комментариев и эссе об иудаизме (в том числе ''Трудная свобода'' (1963) и ''Девять талмудических чтений'' (1968 / 1977)).}. Влияние феноменологии обусловлено в том числе возможностью понимания почти все последующих теоретических формирований в немецкой и францезской философии либо как расширение феноменологии, либо как реакцию на нее. Понимание феноменологии, значит, необходимо для понимания последующего развития философии XX века.

За прошедшие годы феноменология сделала значительный вклад во многие области философии и предложила новаторский анализ интенциональности, восприятия, воплощения, эмоций, самосознания, интерсубъективности, темпоральности, историчности и истины. Она целенаправлено критикует редукционизм, объективизм и сциентизма. Подробно рассматривая человеческое существование и понимая субъекта как воплощенное и социально и культурно встроенное бытие-в-мире, феноменология делает решающий вклад в такие эмпирические дисциплины как психиатрия, социология, психология, литературоведение, антропология и архитектура.

Многие из ныне классических работ по феноменологии были написаны в первой половине XX века, но в последние годы интерес к ней возобновился. Феноменология сейчас переживает своего рода ренессанс.

Феноменология развивалась как во многом разнородное движение с множеством ответвлений, все феноменологи после Гуссерля дистанцировались от различных аспектов его исходной программы, и вряд ли корректно говорить о феноменологии как о философской системе с четко очерченным корпусом доктрин, однако не следует упускать из виду те всеобъемлющие проблемы и темы, которые объединяли и продолжают объединять ее сторонников.

В введении, конечно, невозможно отдать должное всем феноменологам, и я буду опираться прежде всего на работы Гуссерля, раннего Хайдеггера и Мерло-Понти --- трех мыслителей, чье влияние на развитие феноменологии бесспорно. Многие вводные книги по феноменологии содержат отдельные главы, посвященные этим мыслителям, но я поступлю иначе: вместо указания на их различия --- часто, на мой взгляд, преувеличивающиеся из-за фундаментально ошибочного толкования идей Гуссерля --- я сосредоточусть на том, что в них общего.

В первой части книги рассмотрена концепция философии согласно феноменологии --- вопрос метода, фокус на виде от первого лица, анализ проживаемого мира --- также развитие феноменологической традиции.

Во второй части я откладываю в сторону методологические соображения чтобы рассмотреть несколько примеров и моделей феноменологического анализа. Я рассмотрю феноменологические исследования пространственности и воплощенности, а также анализ интерсубъективности и сообщества.

Третья часть книги демонстрирует применение феноменологии вне философии. В частности, я опишу ее влияние на социологию, психологию и когнитивистику.

Главная цель этой книги --- не защита феноменологии, но ее представление: изложение ряда ее идей в максимально доступной форме без защиты их от разного рода критики. Я старался как можно меньше внимания уделять интерпретационным разногласиям и противоречиям. Так что если вы ищете систематическую защиту феноменологии и широкое соотнесение ее положений и выводов с исследовательской литературой, эта книга вам не подойдет\footnote{См., например, Захави 2003, 2005, 2014, 2017 и Галлахер и Захави 2012.}.

\end{document}
